% Generated from validated Article Draft Result. Do not hand-edit; revise the structured draft instead.

3月25日のGPT-4o image generationは、画像生成をGPT-4oのnative capabilityとしてChatGPTの対話フローへ統合した。画像生成が独立した専門モデルUIだけでなく、general-purpose multimodal interactionの一部へ入った出来事である。\autocite{src-877e7da88408}


3月27日のQwen2.5-Omniはtext、image、audio、videoを入力として扱い、streaming speech outputまで含むend-to-end multimodal modelとして公開された。入出力modalitiesの組合せ自体がmodel architectureとUXの設計対象になった。\autocite{src-de4ffd746276}


5月のGoogle I/OではVeo 3、Imagen 4、Flow、Lyria 2が同時期に示され、video、image、filmmaking workflow、musicまでgenerative media surfaceが広がった。個々のmodel identityを保ちつつ、制作環境としての束を読む必要がある。\autocite{src-e68f8ce96685}


multimodalという言葉には、複数modalitiesを理解する能力と、複数mediaを生成する能力が混在する。H1では後者が急速に前景化し、model evaluationだけでなく編集・制作workflowとの接続が重要になった。\autocite{src-877e7da88408,src-de4ffd746276,src-e68f8ce96685}


ただしmedia generationは、そのまま外部世界で行動するagent actionと同義ではない。画像・動画・音声の生成surface拡大と、browser/tool/sandboxを操作するexecution surfaceは、相互に近づきつつも別の技術層として区別する。\autocite{src-877e7da88408,src-e68f8ce96685}


\begin{claimboundary}
一次資料から確認できるmodalities、release scope、product integrationを記録する。画質・音質・動画品質などの優位性は、vendor評価を独立再現した結果としては扱わない。
\end{claimboundary}
